\section{Background}
\label{sec:background}

\todo{Synthesize from docs/theoretical-foundations.md. Cover the
theoretical foundations that ground the taxonomy.}

\subsection{Cognitive Neuroscience of Memory}

\subsubsection{Complementary Learning Systems}
McClelland~\etal~\cite{mcclelland1995} proposed that the brain uses two
complementary systems to avoid catastrophic forgetting: a fast learner
(hippocampus) that records raw episodes and a slow learner (neocortex)
that extracts stable patterns through interleaved replay.

\subsubsection{Hebbian Learning}
``Neurons that fire together wire together''~\cite{hebb1949}. Long-term
potentiation (LTP) and long-term depression (LTD) provide the
biological basis for associative memory~\cite{bliss1973}.

\subsubsection{Spreading Activation}
Collins and Loftus~\cite{collins1975} showed that activating a concept
propagates through a weighted semantic network, priming associated
concepts. This explains how ``doctor'' activates ``nurse'' even without
direct keyword overlap.

\subsubsection{Bjork Dual-Strength Forgetting}
Each memory has two independent strengths: \emph{storage strength} (how
well-encoded) and \emph{retrieval strength} (how accessible
now)~\cite{bjork1992}. Retrieval strength decays over time; storage
strength only increases with access. This model explains tip-of-the-tongue
phenomena and why relearning is faster than initial
learning. Anderson~\etal~\cite{anderson1994} showed that retrieving a
memory suppresses competing memories sharing the same cues
(retrieval-induced forgetting, RIF).

\subsection{Information Retrieval}

\subsubsection{Reciprocal Rank Fusion}
RRF merges multiple ranked result sets without score
calibration~\cite{cormack2009rrf}: $\text{score}(d) = \sum_i 1 / (k +
\text{rank}_i + 1)$. Bruch~\etal~\cite{bruch2023} showed that convex
combination outperforms RRF when tuning data is available, but RRF is
superior in zero-shot settings---the typical case for agent memory.

\subsection{Yogac\=ara Buddhist Psychology}

\todo{Expand this section. This is a key differentiator of the paper.
Draw from docs/theoretical-foundations.md.}

The Yogac\=ara school of Buddhist philosophy, systematized in the
\emph{Yogac\=arabh\=umi-\\'s\=astra} attributed to
Maitreya/Asa\.{n}ga and the \emph{Cheng Weishi Lun} (\cjk{成唯識論})
by Xuanzang (659~CE), proposes eight consciousnesses. The eighth,
\emph{\=alaya-vij\~n\=ana} (\cjk{阿賴耶識}, storehouse consciousness),
is a persistent substrate holding all \emph{b\=ija} (seeds)---latent
potentials that ripen into experience when conditions
converge~\cite{yogacarabhumi, chengweishilun}. Seeds are perfumed
(\emph{v\=asan\=a}, \cjk{薰習}) by each experience, accumulating
traces that shape behavior. Periodic transformation
(\emph{\=a\'sraya-par\=av\d{r}tti}, \cjk{轉依}) restructures the
storehouse itself.
